\section{The objections that could break the argument}
\label{sec:objections}

\subsection{Reliability and distribution shift}

An 80\% task horizon is not 99.9\% gate reliability. Even zero failures in a trial cannot establish
universal safety. With $n$ independent, representative Bernoulli trials and no failures, the one-sided
95\% binomial upper bound is $1-0.05^{1/n}$. Reaching $10^{-3}$ needs 2,995 such trials; reaching
$5\times10^{-5}$ needs 59,914. Repeated prompts on the same few dossiers do not create that many
independent cases. Volume is the practical difficulty: a route with five visits a month yields sixty
visits a year, and sixty visits without failure bound the error rate only below 4.9\%. A year without
incident on a low-volume route therefore cannot support a 0.1\% ceiling. The qualification plan must
say where the volume comes from, off-production dossiers, dedicated campaigns, or several sites, and
why those cases are representative and independent; sharing a service across customers helps to
accumulate observations without making their contexts identical. Changes in policy, supplier population, or tools can invalidate transport of the
estimate.

A defensible system is qualified on declared domains and monitored for drift; unknown cases trigger a
valid recovery, inquiry, or refusal. If the accepted contract requires a human recovery, the
replacement definition has not been met. If the system merely refuses an expanding fraction of
legitimate work, its false-refusal and utility metrics fail. Neither a high average score nor an
unlimited abstention policy closes a gate.

Reliability is also a path property. Let $\mathcal{C}_i$ be correctness of gate $i$ under an
independent reference standard. The probability of no gate errors is exactly
\begin{equation}
\Pr\Bigl(\bigcap_{i=1}^{n}\mathcal{C}_i\Bigr)=\Pr(\mathcal{C}_1)\prod_{i=2}^{n}
\Pr\Bigl(\mathcal{C}_i \,\Big|\, \bigcap_{j<i}\mathcal{C}_j\Bigr), \label{eq:chain}
\end{equation}
when the conditioning events have positive probability. Twenty marginal accuracies of 0.99 imply a
no-error probability of $0.99^{20}\simeq 0.818$ under independence, but only the Fr\'echet bounds
0.80--0.99 without a dependence model. Two agents using the same model, policy interpretation, or false
source may agree and fail together. Common failures can have worse severity despite a higher
probability of an entirely correct run. Route tests must measure dependence and loss, not infer
diversity from the number of agent instances. Published evaluations cover pieces of this problem:
enterprise web tasks \citep{drouin2024}, operational IT scenarios \citep{jha2025}, long software-task
completion \citep{kwa2025}, and contextual integrity in enterprise workflows \citep{fu2026}. They
motivate separate utility, reliability, and information-flow evaluation; they do not supply the DGF
scenario parameters, and they establish present limits, not permanent ones.

\subsection{The adversary and the project seeking approval}
\label{sec:strategic}

A supplier document is also an attack surface. AgentDojo evaluates how untrusted tool-returned data can
redirect agents toward malicious actions \citep{debenedetti2024}. A DGF agent reads contracts, tickets,
and inventories supplied by parties with interests in its disposition. Instruction filtering alone is
not an assurance argument; external text must not gain the authority to change rules or grant tools.

Projects can also optimize the measured criteria without any injection. The allocation of effort across
differently rewarded activities connects to multitask incentives \citep{holmstrom1991}; the distinction
between improvement and gaming is explicit in \citet{kleinberg2019}; strategic classification supplies
the broader setting in which evaluated people change the features presented to a decision rule
\citep{hardt2016}. Consider a minimal gate model. A project begins with true compliance $\xi$. It
chooses substantive remediation $u\ge 0$ and cosmetic reporting effort $g\ge 0$, producing observable
score $y=\xi+u+g+\varepsilon$. A threshold $\tau$ determines preliminary approval. Let $V$ be the
project value of that approval, $c_u(u)$ and $c_g(g)$ the two effort costs, and $D(g,a)$ the expected
sanction, lost approval value, and remediation cost from detecting misleading reporting under assurance
intensity $a$. The project chooses
\begin{equation}
\max_{u,g\ge 0}\; V\Pr(y\ge\tau)-c_u(u)-c_g(g)-D(g,a). \label{eq:strategic}
\end{equation}
Let $p(s)=\Pr(s+\varepsilon\ge\tau)$ and $s=\xi+u+g$. At a differentiable interior optimum,
$Vp'(s)=c_u'(u)=c_g'(g)+D_g(g,a)$. Assurance raises the marginal cost of gaming when $D_{ga}>0$. With
$\omega_u=c_u''>0$, $\omega_g=c_g''+D_{gg}>0$, $t=Vp''(s)\le 0$, and
$\Delta=\omega_u\omega_g-t(\omega_u+\omega_g)>0$, implicit differentiation gives
\begin{equation}
\frac{\partial u^{*}}{\partial a}=-\frac{tD_{ga}}{\Delta}\ge 0,\qquad
\frac{\partial g^{*}}{\partial a}=\frac{(t-\omega_u)D_{ga}}{\Delta}<0,\qquad
\frac{\partial s^{*}}{\partial a}=-\frac{\omega_uD_{ga}}{\Delta}<0 . \label{eq:statics}
\end{equation}
Stronger assurance reduces cosmetic effort and, with diminishing marginal approval returns, weakly
increases genuine remediation. The observable score still falls, so a lower pass rate after assurance
is strengthened is not by itself evidence of worse projects. For smooth zero-centered symmetric noise
with density $f$ strictly decreasing away from its mode, $p''(s)=-f'(\tau-s)$; hence $t\le 0$ when
$s\ge\tau$, where preliminary passage probability is at least one half. Below the threshold, where
$t>0$, both efforts can fall. This predicts heterogeneous responses by pre-intervention distance to the
threshold. Flat density segments and corner solutions require separate treatment; clearer criteria can
lower both effort costs, leaving their allocation effect ambiguous. The ancillary mathematical notes
give the derivation and boundary checks.

The model identifies a testable failure mode: declining escalations can coexist with concealed defects.
A falling exception rate can reflect improved systems, better extraction, or concealed exceptions.
Readiness inferred only from submitted documentation becomes less informative as projects learn the
gate. An enterprise test compares claimed and independently verified compliance, uses hidden holdout
checks of stable requirements, and records post-approval defects. Provenance checks, random audits,
independent validation, and restricted authority all enter the assurance cost. The law requires
automation that remains effective against adaptation, not a faster route to approving misleading
evidence.

\subsection{Who verifies the automated verifier?}

The assurance argument is a division of mechanisms, not an infinite stack of identical language-model
opinions. Typed outputs and capability-limited workers enforce action constraints. Independent
retrieval checks provenance against authoritative systems. Deterministic tests establish narrow
invariants. Differently constructed verifiers and held-out adversarial tests address correlated
interpretation errors. Recorded outcomes expose deterioration over time. Model diversity only helps
when the relevant errors actually differ; a shared false policy can defeat every model. An agent must
not be able to create its own policy exception, alter the evidence against which its output is tested,
or confer on itself a privileged tool.

A safety case then states the accepted residual risk, domains, fallback, revocation conditions,
evidence of performance, and accountable operator. Vendor responsibility, contractual remedies, and
insurance can allocate loss and create incentives; they do not establish correctness or independence.
A mandatory human audit or intervention remains human work in $B_A$ or $h_E$. Bainbridge's analysis
explains why supervisory and recovery costs cannot be ignored and argues for continuing human
involvement \citep[pp.~775--777]{bainbridge1983}. The law of gate substitution challenges the permanence of that
residual, not the existence of those costs.

The strong thesis predicts that this recurring assurance work itself becomes executable under accepted
rules. It can fail: new policies or distribution shifts may continually require human judgments that
no validated procedure replaces. A standing mandate solves the authority of execution only within its
scope; it does not solve the epistemology of assurance or discharge statutory management duties.

\subsection{Why not repeat profession-wide automation predictions?}

A diagnostic classification is not a whole radiology practice, and vehicle control is not every
function in road transport. The same warning applies here: passing a checklist is not the whole DGF.
The advantage of the gate unit is that it enumerates the inquiry, output, authority, and support that
remain before claiming replacement. It also provides a finite case registry against which failure can
be observed.

The informational character of the work is an advantage, not a metaphysical guarantee. Many DGF inputs
and actions are already digital; others require new observations, negotiation, or physical inspection.
A human collecting indispensable evidence outside the buyer's boundary still performs necessary work.
Appendix~\ref{app:limits} makes the information limitation precise. The Integrate route is included because inherited evidence and low repetition test this weakness rather than allowing the thesis to be demonstrated only
on a clean software pipeline.

\subsection{A difficult residual is not a permanently human residual}

Knowledge hierarchies concentrate expert work on problems that other workers cannot resolve
\citep[Sections~1 and~3]{garicano2015}. Automation can initially reproduce that arrangement: machines
handle standard cases and people absorb exceptions. The law does not confuse this intermediate
allocation with its endpoint. The exception handler has its own observations, hypotheses, tests, and
outputs; it belongs to the substitution target. An undocumented acquisition incident may require fresh
evidence, not better text generation. Supplier negotiations require interaction with another party. A
new project's failure can reveal an unknown architecture rather than a missing checklist field. The
fact that a current system cannot close these activities is evidence about the remaining technical
work, not a proof that the role must always be human. Hypothesis~H6 turns this objection into a
measurement.

\subsection{Exhaustion or renewal: the rival dynamics}
\label{sec:renewal}

Continual automation is not enough to reach zero. Let $R_t$ be the residual human load, $a$ the
fraction of it closed in each period, and $b$ the new human needs that appear in each period, some of
them created by the automation itself:
\begin{equation}
R_{t+1}=(1-a)R_t+b,\qquad 0<a<1,\qquad R_\infty=\frac{b}{a}. \label{eq:renewal}
\end{equation}
With constant $b>0$, automation keeps progressing and the human requirement never disappears. The law
therefore needs more than $a>0$: it needs $b_t\to 0$, new indispensable human work that stops
appearing. Convergence is also weaker than closure: a load that tends to zero can still be positive at
any finite date, as Section~\ref{sec:workforce} notes for posts. Such a residual is compatible
with the separate 80\% headcount reduction predicted for 2033.

The reason to expect exhaustion rather than renewal is that the work created by automation has the
same shape as the work it replaces. Supervising, evaluating, triaging exceptions, and maintaining
agents all read records and return artifacts, so the same substitution applies to them, one level up.
If each unit of automated work creates $s<1$ units of supervisory work and each level of supervision
can in turn be executed under mandate, the human need that remains after $k$ levels is of order $s^k$,
not a constant. The rival is a level at which $s\ge 1$, or at which no mandate is accepted, for
instance a statutory duty of the management body. This is an argument, not evidence, and it is an argument about contraction, not about closure: it
explains why the residual can become small. Closure is the architectural condition of
Section~\ref{sec:conditions}, an implementation in which no necessary human step remains, including
for its maintenance, and it is established contract by contract, not as the limit of a sequence.
Hypothesis~H6 and the intervention registry of Section~\ref{sec:testing} are how the contraction is
checked.

\subsection{The economic rival: a small but valuable human residual}

The strongest rival is not that inputs and outputs are absent. \citet{gans2026} emphasize O-ring task
complementarity: improving automated tasks can increase the importance of remaining manual bottlenecks
and change the incentive to automate them. This gives a serious rival endpoint: a small but valuable
human residual may remain economically preferable. The law of gate substitution needs the quality-adjusted machine
frontier eventually to overtake that residual too. It cannot derive that crossing from the current
size of a department. Section~\ref{sec:positioning} places this rival among the closest economic models.
